%%%%%%%%%%%%%%%%%%%%%%%%%%%%%%%%%%%%%%%%%%%%%%%%%%%%%%%%%%%%%%%%%%%%%%
%
% Institut für Information Systems Engineering
% Forschungsgruppe INSO
% Arbeitsgruppe ESSE
% https://establishing-security.at/
% lva.security@inso-world.com
%
%%%%%%%%%%%%%%%%%%%%%%%%%%%%%%%%%%%%%%%%%%%%%%%%%%%%%%%%%%%%%%%%%%%%%%

\documentclass[12pt,a4paper,titlepage,oneside]{scrartcl}
\newcommand{\lang}{de}
\usepackage{esseProtocol}

%%%%%%%%%%%%%%%%%%%%%%%%%%%%%%%%%%%%%%%%%%%%%%%%%%%%%%%%%%%%%%%%%%%%%%
%
% FOR STUDENTS
%
%%%%%%%%%%%%%%%%%%%%%%%%%%%%%%%%%%%%%%%%%%%%%%%%%%%%%%%%%%%%%%%%%%%%%%

% Team number or "0" for Lab0
%TODO team number
\newcommand{\team}{44}
% Date
%TODO fill in creation date
\newcommand{\datum}{\today}
%TODO lab number
% valid values: "Lab0", "Lab1" (be sure to use Uppercase for first character)
\newcommand{\lab}{Lab1}

%TODO name of course
\newcommand{\lvaname}{Einführung in Security} % if 6 ECTS variant, otherwise: Introduction to Security
%TODO number of course
\newcommand{\lvanr}{194.157} % if 6 ECTS variant, otherwise: 183.594
%TODO year and term, for example: "2024 S" or "2024 W", etc.
\newcommand{\semester}{2024 W}

% Student data in Lab0 or 1. student of team in Lab1
\newcommand{\studentAName}{Kevin Csele}
\renewcommand{\studentAMatrnr}{12122544}

% 2. student of team in Lab1, for Lab0 or if your team has less students, remove these 2 lines
\newcommand{\studentBName}{Clemens Schneider}
\renewcommand{\studentBMatrnr}{MATRIKELNUMMER}

% 3. student of team in Lab1, for Lab0 or if your team has less students, remove these 2 lines
\newcommand{\studentCName}{Luka Twaroch}
\renewcommand{\studentCMatrnr}{MATRIKELNUMMER}

% 4. student of team in Lab1, for Lab0 or if your team has less students, remove these 2 lines
\newcommand{\studentDName}{Wen Long Zhou}
\renewcommand{\studentDMatrnr}{MATRIKELNUMMER}

% 5. student of team in Lab1, for Lab0 or if your team has less students, remove these 2 lines
\newcommand{\studentEName}{Ramin Shaikh}
\renewcommand{\studentEMatrnr}{12123657}

%%%%%%%%%%%%%%%%%%%%%%%%%%%%%%%%%%%%%%%%%%%%%%%%%%%%%%%%%%%%%%%%%%%%%%
%
% DO NOT CHANGE THE FOLLOWING PART
%
%%%%%%%%%%%%%%%%%%%%%%%%%%%%%%%%%%%%%%%%%%%%%%%%%%%%%%%%%%%%%%%%%%%%%%

\newcommand{\colormode}{color}
\newcommand{\dokumenttyp}{Abgabedokument \lab}

\begin{document}

\maketitle
\setcounter{section}{0}
\setcounter{tocdepth}{2}
\tableofcontents

%%%%%%%%%%%%%%%%%%%%%%%%%%%%%%%%%%%%%%%%%%%%%%%%%%%%%%%%%%%%%%%%%%%%%%
%
% CONTENT OF DOCUMENT STARTS HERE
%
%%%%%%%%%%%%%%%%%%%%%%%%%%%%%%%%%%%%%%%%%%%%%%%%%%%%%%%%%%%%%%%%%%%%%%
\newpage

\section{Der Service war auch schon besser ...}

\subsection{Achtung! Streng geheim!}
Um diese Aufgabe zu l\"osen, hat es gen\"ugt, das besagte PDF im Browser zu \"offnen. Der ``streng geheime`` String befand sich im Titel des Tabs.

\subsection{Eine schräge Nummer}
Die Rechnungsnummer wurde zwar von einem schwarzen Rechteck verdeckt, ließ sich jedoch ganz einfach kopieren, indem man die betroffene Stelle markiert -> Strg + C

\subsection{Was letzte Preis?}
Selbes Spiel, auch der Preis ließ sich ganz simpel herauskopieren.

\subsection{IBANs sollte man verbannen!}
Um den IBAN aufzudecken, habe ich PDF-XChange verwendet, um das schwarze Rechteck mit dem Objektbearbeitungswerkzeug zu entfernen.


\section{Wireless Time Travel}

\subsection{Vier zukunftssichere Handschl\"age}
Nicht gelöst.

\subsection{Code der Zukunft}
Nicht gelöst.

\subsection{Ungew\"ohnlich verschl\"usselte Botschaft}
Nicht gelöst.

\subsection{Geheimnisvoller Zugang: superboss}
Nicht gelöst.

\subsection{Unbrauchbarer Schl\"ussel}
Nicht gelöst.

\subsection{Einen Schl\"ussel f\"ur einen Schl\"ussel! Echt jetzt?!}
Nicht gelöst.

\subsection{Verborgenes Protokoll}
Nicht gelöst.


\section{Bot Bot Bot Bot}

\subsection{I keep you my little secret ...}
Nicht gelöst.


\section{C\"asars Schl\"usselbund}

\subsection{Schl\"ussel. Knacken.}
Nicht gelöst.

\subsection{Passw\"orter Retten.}
Nicht gelöst.


\section{Paranoider Mozart}

\subsection{MozART.}
Nicht gelöst.


\section{Zertifiziertes Durcheinander}

\subsection{Zertifizieren ist schwer}
Um den Certificate Signing Request zu erstellen habe ich den folgenden Befehl verwendet:
\lstinline{openssl req -newkey rsa:4096 -sha512 -config openssl.cnf -out csr.csr -subj "/CN=12123657-Intermediate-CA-WS2024/OU=ESSE-Lab1-Exercise"} \\
\lstinline{req} ist der Befehl um einen CSR zu erstellen; \\
\lstinline{-newkey rsa:4096} spezifiziert, dass ein neuer Key (4096-Bit RSA) erstellt werden soll; \\
\lstinline{-sha512} gibt an, dass sha512WithRSAEncription als Signaturalgorithmus verwendet werden soll; \\
mit \lstinline{-config} wird angegeben, welches config file zu verwenden ist; \\
\lstinline{-out} bestimmt das output-file und \\
\lstinline{-subj "/CN=12123657-Intermediate-CA-WS2024/OU=ESSE-Lab1-Exercise"} definiert die gewünschten Namens-Parameter im CSR. \\ \\
Das config file dient dazu, die nötigen X509v3 Parameter zu setzen und sieht aus wie folgt:
\begin{lstlisting}[caption=openssl.cnf,label=code:opensslConfig,style=simple]
  [ req ]
  default_bits           = 4096
  default_md             = sha512
  default_keyfile        = privkey.pem
  distinguished_name     = req_distinguished_name
  req_extensions          = v3_req

  [ req_distinguished_name ]

  [ v3_req ]
  subjectKeyIdentifier = hash
  basicConstraints = critical, CA:true, pathlen:0
  keyUsage = critical, Certificate Sign, CRL Sign
\end{lstlisting}
Nach Ausführung des oben genannten Befehls, wird der CSR in der Datei csr.csr gespeichert, diese wurde im Abgabetool eingereicht.


\section{Zeitreise durch das World Wide Web}

\subsection{Wieder Elvis}
Nicht gelöst.

\subsection{C\"asarMussWeg! MussC\"asarWeg?}
Nicht gelöst.

\subsection{dackboor.}
Nicht gelöst.

\subsection{Schlechtes Timing (Time Travel Edition)}
Nicht gelöst.

\subsection{Sorcerer ... ?}
Nicht gelöst.


\section{Seitlich fließend}

\subsection{Newton und Co KG.}
Nicht gelöst.


\section{Antike Mobile Security}

\subsection{iTimeTravel}
Nicht gelöst.

\subsection{AND(roid)ERS}
Nicht gelöst.


\section{Babycam Espionage}

\subsection{The Rise of the HuManoiD5}
Nicht gelöst.

\subsection{ETA}
Nicht gelöst.


\section{Das Social Media der Zukunft}

\subsection{Der vergiftete Passwort Reset}
Nicht gelöst.

\subsection{Accountübernahme}
Nicht gelöst.


\section{Hidden Timelines}

\subsection{Phantom Domain}
Nicht gelöst.


\section{Vault Voyage}

\subsection{That's all your vault!}
Nicht gelöst.


\section{Wikinger Overflow}

\subsection{\"Uberlauf. Hand drauf.}
Nicht gelöst.

\subsection{Typisch Typing ... Stufe 1}
Nicht gelöst.

\subsection{Typisch Typing ... Stufe 2}
Nicht gelöst.

\subsection{Typisch Typing ... Stufe 3}
Nicht gelöst.


\section{Tap to the Future}

\subsection{Tick Tock Tap}
Nicht gelöst.


\section{So viele}

\subsection{Das Device ist heiß}
Nicht gelöst.

\subsection{Persona non grata}
Nicht gelöst.

\subsection{Eine Frage der Kommunikation}
Nicht gelöst.

\subsection{Treffpunkt}
Nicht gelöst.

\subsection{Alles dokumentiert!}
Nicht gelöst.

\subsection{Es geht immer um Inhalte}
Nicht gelöst.


\section{Web of Treats}
Nicht gelöst.

\subsection{Mitgliedschaftsnr.}
Nicht gelöst.

\subsection{Geheimer Artikel}
Nicht gelöst.

\subsection{\"Uberf\"ullt}
Nicht gelöst.

\subsection{A shell in the forest?}
Nicht gelöst.

\subsection{Elvis}
Nicht gelöst.


\section{Das. Beste. Text. Adventure. Aller. Zeiten.}

\subsection{Time to travel!}
Nicht gelöst.

\subsection{Mein Name?}
Nicht gelöst.

\subsection{Ein PIN!}
Nicht gelöst.

\subsection{Ach ... ein Schl\"ussel}
Nicht gelöst.

\subsection{Flag!}
Nicht gelöst.


\section{Passwörter werden wir auch nie los, oder?!}

\subsection{Gute Idee, um ein Passwort zu verstecken?!}
Nicht gelöst.

\subsection{Call Julius ... äh. John.}
Nicht gelöst.

\subsection{Nicht nur Ziffern, sonder auch ...?}
Nicht gelöst.

\subsection{/etc/ANTIK?}
Nicht gelöst.

\subsection{Sicher sicher?}
Nicht gelöst.

\subsection{Zeitlose Liste}
Nicht gelöst.

\subsection{(Image)magic(k)}
Nicht gelöst.

\subsection{Auch in Zukunft ein schweres Passwort?}
Nicht gelöst.


\section{Franz Joseph und die Kommandozeile}

\subsection{Stage}
Nicht gelöst.

\subsection{Stagee}
Nicht gelöst.

\subsection{Stageee}
Nicht gelöst.

\subsection{Stageeee}
Nicht gelöst.

\subsection{Stageeeee}
Nicht gelöst.

\subsection{Stageeeeee}
Nicht gelöst.

\subsection{Stageeeeeee}
Nicht gelöst.

\subsection{Stageeeeeeee}
Nicht gelöst.

\subsection{Stageeeeeeeee}
Nicht gelöst.

\subsection{Stageeeeeeeeee}
Nicht gelöst.

\subsection{Stageeeeeeeeeee}
Nicht gelöst.

\subsection{Stageeeeeeeeeeee}
Nicht gelöst.


\section{Ueberschrift 1}

\subsection{Hinweise}
\emph{Hinweise:}
\begin{itemize}
    \item Verwenden Sie entweder diese deutsche Version oder die englische Version in \lstinline{protocol.tex}.
    \item Setzen Sie alle Variablen nach \emph{FOR STUDENTS} in der .tex Datei.
    \item Ersetzen Sie die Platzhalter für Ihre Namen und MatNr.
    \item Löschen Sie diese Sektion über Hinweise und die folgenden Beispiel-Kapitel.
    \item Achten Sie auf geforderte Formate und Anforderungen an die Dateinamen.
    \item Führen Sie \lstinline{pdflatex} mindestens zweimal aus, damit die Referenzen und Seitenzahlen richtig im PDF dargestellt werden.
    \item Sie können dazu auch das Makefile verwenden: \lstinline{make de}.
\end{itemize}

\section{Beispiele}

\subsection{Source Code formatieren}
Es folgen einige Beispiele wie Sourcecode in diesem Dokument formatiert und referenziert werden kann
(\hyperref[code:beispiel1]{siehe Listing~\ref*{code:beispiel1} auf Seite~\pageref*{code:beispiel1}} und \hyperref[code:beispiel2]{siehe Listing~\ref*{code:beispiel2} auf Seite~\pageref*{code:beispiel2}}).

Ebenso können kurzer Code oder kurze Befehle direkt in der Zeile in einem \lstinline{lstinline Block} mit typengleicher Schrift formatiert werden.

\lstinputlisting[caption=Example C/C++ file,label=code:beispiel1,style=c]{example.c}

\begin{lstlisting}[caption=Example bash script,label=code:beispiel2,style=simple]
#!/bin/bash
echo "Bash version ${BASH_VERSION}..."
for i in {0..10..2}
  do
     echo "Welcome $i times"
 done

echo "some very very very very very very very very very very very very very very very very very very very very long string"

exit 0;
\end{lstlisting}

\subsection{Bilder}

Es folgen einige Beispiele wie Bilder in diesem Dokument eingefuegt werden koennen
(\hyperref[fig:logo1]{siehe Abbildung~\ref*{fig:logo1} auf Seite~\pageref*{fig:logo1}}).

\begin{figure}[h!]
  \centering
  \fbox{
    \includegraphics[width=0.4\textwidth]{./imgs/logos/esse-logo-color.png}
  }
  \caption{ESSE Logo}
  \label{fig:logo1}
\end{figure}


%%%%%%%%%%%%%%%%%%%%%%%%%%%%%%%%%%%%%%%%%%%%%%%%%%%%%%%%%%%%%%%%%%%%%%
%
% DO NOT CHANGE THE FOLLOWING PART
%
%%%%%%%%%%%%%%%%%%%%%%%%%%%%%%%%%%%%%%%%%%%%%%%%%%%%%%%%%%%%%%%%%%%%%%

\end{document}